\documentclass[11pt]{article}
\usepackage[margin=1in]{geometry}
\usepackage{amsmath,amssymb,amsthm,mathtools}
\usepackage{graphicx}
\usepackage{hyperref}
\usepackage{microtype}

\hypersetup{colorlinks=true,linkcolor=blue,citecolor=blue,urlcolor=blue}

\newtheorem{theorem}{Theorem}
\newtheorem{lemma}{Lemma}
\newcommand{\Nstar}{\lVert\cdot\rVert_*}
\newcommand{\ellone}{\ell_1}
\DeclareMathOperator{\domreg}{dom}

\title{A Strictly Convex Counterexample for Distance-Dominance Regions}
\author{Candidate counterexample packet for arXiv:1303.4733}
\date{9 August 2026}

\begin{document}
\maketitle

\begin{center}
\fbox{\parbox{0.91\linewidth}{\textbf{Status.} Candidate counterexample, likely valid, pending independent human review. The construction answers the strict-convexity conjecture in Section 5 of Reem's paper.}}
\end{center}

\section{Source problem}

Daniel Reem's paper \emph{Topological properties of sets represented by an inequality involving distances}, arXiv:1303.4733v2 \cite{Reem}, proves the following for a convex subset of a uniformly convex normed space and nonempty sites $P,A$ with $d(P,A)>0$:
\[
 \partial\domreg(P,A)=\{x:d(x,P)=d(x,A)\},
 \qquad
 \operatorname{int}\domreg(P,A)=\{x:d(x,P)<d(x,A)\},
\]
and $\domreg(P,A)$ is the closure of the strict-inequality set, where
\[
 \domreg(P,A):=\{x:d(x,P)\le d(x,A)\}.
\]
In Section 5 (PDF page 11), the paper asks whether uniform convexity can be weakened to strict convexity and states: ``We conjecture that in this case there are counterexamples to Theorem 3.3.'' The source passage is reproduced below.

\begin{figure}[ht]
  \centering
  \includegraphics[width=0.90\textwidth]{figures/open_problem_crop.png}
  \caption{The strict-convexity conjecture in Section 5 of arXiv:1303.4733v2, PDF page 11.}
\end{figure}

\section{Proof intuition}

The ordinary $\ell_1$ geometry makes the desired dominance relation transparent: a vector supported on two coordinates can be at least as close as a basis vector because it can absorb positive mass from either coordinate. Ordinary $\ell_1$ is not strictly convex. We make it strictly convex by adding a Hilbertian term whose coordinate weights tend to zero. The perturbation enforces strict convexity globally but becomes invisible along the tail.

Take the first site to be all basis vectors and the second site to be all averages of two distinct basis vectors. These sites remain one unit apart in the $\ell_1$ part of the norm. Their distances from the origin both have infimum one, attained only asymptotically. For every point in a fixed ball, a second averaging coordinate can be sent far enough into the tail that the Hilbertian cost vanishes; the remaining coordinate then gives a strictly smaller Hilbertian contribution than the corresponding basis vector. Thus the whole ball lies in the dominance region although the origin is an equality point.

\section{Counterexample}

Let $X=\ell_1(\mathbb N;\mathbb R)$, put $w_n=2^{-n}$, and define
\[
 q(x):=\left(\sum_{n=1}^{\infty}w_nx_n^2\right)^{1/2},
 \qquad
 \lVert x\rVert_*:=\lVert x\rVert_1+q(x).
\]
Write $(e_n)$ for the standard unit vectors and set
\[
 A:=\{e_n:n\in\mathbb N\},
 \qquad
 P:=\left\{\frac{e_n+e_m}{2}:n,m\in\mathbb N,\ n\ne m\right\}.
\]

\begin{theorem}
The norm $\Nstar$ makes $X$ a strictly convex Banach space. The sets $P$ and $A$ are closed and satisfy
\[
 d_*(P,A)=1,
 \qquad
 d_*(0,P)=d_*(0,A)=1.
\]
Moreover,
\[
 B_*(0,1/2)\subseteq \domreg(P,A).
\]
Consequently, $0$ is an interior point of $\domreg(P,A)$ at which $d_*(0,P)=d_*(0,A)$. Hence
\[
 \partial\domreg(P,A)\ne\{x:d_*(x,P)=d_*(x,A)\}
\]
and
\[
 \operatorname{int}\domreg(P,A)\ne\{x:d_*(x,P)<d_*(x,A)\}.
\]
In particular, strict convexity cannot replace uniform convexity in Theorem 3.3 of \cite{Reem} for arbitrary positively separated sites.
\end{theorem}

\begin{proof}
Since $q(x)\le \lVert x\rVert_2\le \lVert x\rVert_1$,
\[
 \lVert x\rVert_1\le \lVert x\rVert_*\le 2\lVert x\rVert_1.
\]
Thus $\Nstar$ is an equivalent Banach norm on $\ell_1$. The norm $q$ is a strictly convex norm because every $w_n$ is positive. If equality held in the midpoint inequality for $\Nstar$ at two distinct unit vectors, equality would have to hold for both the $\ell_1$ and $q$ summands. Equality for $q$ forces the two vectors to be nonnegative scalar multiples of one another, and equality of their $\Nstar$-norms then forces them to coincide. Hence $\Nstar$ is strictly convex.

For $p_{n,m}:=(e_n+e_m)/2$ and $e_k\in A$,
\[
 \lVert p_{n,m}-e_k\rVert_1=
 \begin{cases}
 1,&k\in\{n,m\},\\
 2,&k\notin\{n,m\}.
 \end{cases}
\]
Therefore $d_*(P,A)\ge1$. Sending two distinct indices to infinity gives
\[
 \lVert p_{n,m}-e_n\rVert_*
 =1+\frac12\sqrt{w_n+w_m}\longrightarrow1,
\]
so $d_*(P,A)=1$. Distinct elements of $A$ are at $\ell_1$-distance $2$, and distinct elements of $P$ (viewed as unordered pairs) are at $\ell_1$-distance at least $1$. Thus both sites are uniformly discrete and hence closed.

At the origin,
\[
 \lVert e_n\rVert_*=1+\sqrt{w_n}\longrightarrow1
\]
and, as $n,m\to\infty$ with $n\ne m$,
\[
 \lVert p_{n,m}\rVert_*=1+\frac12\sqrt{w_n+w_m}\longrightarrow1.
\]
All these vectors have $\ell_1$-norm one, so both infima equal one:
\[
 d_*(0,A)=d_*(0,P)=1.
\]

It remains to prove the ball inclusion. Fix $x\in X$ with $\lVert x\rVert_*<1/2$. Then $|x_j|<1/2$ for every $j$. Let $x_j^+=\max\{x_j,0\}$. For every fixed $n$,
\begin{align}
 \lVert x-e_n\rVert_1
 &=\lVert x\rVert_1+1-2x_n^+,
 \label{eq:l1A}\\
 q(x-e_n)^2
 &=q(x)^2+w_n(1-2x_n).
 \label{eq:qA}
\end{align}
Choose indices $m\to\infty$ with $m\ne n$. Since $x\in\ell_1$, we have $x_m\to0$, and of course $w_m\to0$. The same coordinate calculation gives
\begin{align*}
 \lVert x-p_{n,m}\rVert_1
 &=\lVert x\rVert_1+1-2x_n^+-2x_m^+,\\
 q(x-p_{n,m})^2
 &=q(x)^2+w_n(1/4-x_n)+w_m(1/4-x_m).
\end{align*}
Consequently,
\begin{multline}
 \lim_{m\to\infty}\lVert x-p_{n,m}\rVert_*\\
 =\lVert x\rVert_1+1-2x_n^+
 +\sqrt{q(x)^2+w_n(1/4-x_n)}.
 \label{eq:Plimit}
\end{multline}
Because $x_n<1/2<3/4$,
\[
 q(x)^2+w_n(1/4-x_n)
 <q(x)^2+w_n(1-2x_n).
\]
Equations \eqref{eq:l1A}--\eqref{eq:Plimit} therefore imply
\[
 d_*(x,P)
 \le \lim_{m\to\infty}\lVert x-p_{n,m}\rVert_*
 <\lVert x-e_n\rVert_*
 \qquad(n\in\mathbb N).
\]
In particular, $d_*(x,P)\le\inf_n\lVert x-e_n\rVert_*=d_*(x,A)$. Hence $x\in\domreg(P,A)$, proving
$B_*(0,1/2)\subseteq\domreg(P,A)$.

The origin is therefore an interior point of the dominance region, but its distances to the two sites are equal. It cannot be a boundary point, and both asserted failures follow.
\end{proof}

\section{Verification notes}

The proof uses only coordinate identities. The main points for independent review are:
\begin{enumerate}
  \item positivity of every $w_n$, which makes the Hilbertian summand strictly convex;
  \item the exact $\ell_1$ separation of the two sites;
  \item nonattainment of both distance-one infima at the origin;
  \item the passage $m\to\infty$ in \eqref{eq:Plimit}; and
  \item the radicand gap $w_n(3/4-x_n)>0$ on the stated ball.
\end{enumerate}
A deterministic numerical checker in \texttt{code/verifier.py} evaluates these identities on 500 seeded random finite-support vectors and tests midpoint strictness on a further 500 random pairs. It is a sanity check, not part of the proof.

\section{Novelty and limitations}

This packet gives a counterexample to the strict-convexity conjecture, not a characterization of all strictly convex spaces or all site configurations for which Reem's conclusions remain true. The use of nonattained site distances is essential to this construction and matches the infinite-dimensional case explicitly left outside the auxiliary results discussed by the source paper.

The bounded novelty check performed on 9 August 2026 searched the run registry and attempts for arXiv:1303.4733 and the exact title, and searched arXiv/web indexes with the phrases ``positively separated sites strictly convex Voronoi'', ``fat bisector strictly convex Banach space'', ``uniform convexity to strict convexity Voronoi counterexample'', and the source title together with ``strict convexity''. The search returned the source paper and general finite-dimensional bisector literature, but no later paper explicitly resolving this conjecture and no occurrence of this construction. This is not an exhaustive citation review; novelty confidence is therefore moderate, while mathematical confidence in the self-contained construction is high.

\paragraph{Human-review recommendation.}
Review as a likely full counterexample. The decisive check is the pointwise comparison between \eqref{eq:qA} and \eqref{eq:Plimit}; after that, verify that taking an infimum over $n$ preserves the dominance inequality.

\begin{thebibliography}{9}
\bibitem{Reem}
D. Reem,
\emph{Topological properties of sets represented by an inequality involving distances},
arXiv:1303.4733v2 (2013), Section 5, p.~11.
\url{https://arxiv.org/abs/1303.4733}.
\end{thebibliography}

\end{document}
